\let\negmedspace\undefined
\let\negthickspace\undefined
\documentclass[journal]{IEEEtran}
\usepackage[a5paper, margin=10mm, onecolumn]{geometry}
%\usepackage{lmodern} % Ensure lmodern is loaded for pdflatex
\usepackage{tfrupee} % Include tfrupee package

\setlength{\headheight}{1cm} % Set the height of the header box
\setlength{\headsep}{0mm}     % Set the distance between the header box and the top of the text

\usepackage{gvv-book}
\usepackage{gvv}
\usepackage{cite}
\usepackage{amsmath,amssymb,amsfonts,amsthm}
\usepackage{algorithmic}
\usepackage{graphicx}
\usepackage{textcomp}
\usepackage{xcolor}
\usepackage{txfonts}
\usepackage{listings}
\usepackage{enumitem}
\usepackage{mathtools}
\usepackage{gensymb}
\usepackage{comment}
\usepackage[breaklinks=true]{hyperref}
\usepackage{tkz-euclide} 
\usepackage{listings}
% \usepackage{gvv}                               

\def\inputGnumericTable{}                      
\usepackage[latin1]{inputenc}                    
\usepackage{color}                              
\usepackage{array}                             
\usepackage{longtable}                          
\usepackage{calc}                               
\usepackage{multirow}                           
\usepackage{hhline}                            
\usepackage{ifthen}                          
\usepackage{lscape}
\begin{document}

\bibliographystyle{IEEEtran}
\vspace{3cm}

\title{12.206}
\author{AI25BTECH11024 - Pratyush Panda
}
\maketitle
% \newpage
% \bigskip
{\let\newpage\relax\maketitle}

\renewcommand{\thefigure}{\theenumi}
\renewcommand{\thetable}{\theenumi}
\setlength{\intextsep}{10pt} % Space between text and floats


\numberwithin{equation}{enumi}
\numberwithin{figure}{enumi}
\renewcommand{\thetable}{\theenumi}

\textbf{Question: } \\
The eigenvectors of the matrix 
\begin{align}
\myvec{1 & 2 \\ 0 & 2}
\end{align}
are written in the form $\myvec{1 \\ a}$ and $\myvec{1 \\ b}$. What is $a+b$?
\begin{enumerate}
\item[a] 0
\item[b] $\frac{1}{2}$
\item[c] 1
\item[d] 2
\end{enumerate}
\vspace{0.7cm}

\renewcommand{\theequation}{0.\arabic{equation}}

\textbf{Solution: } \\
Given
\begin{align}
\Vec{A} = \myvec{1 & 2 \\ 0 & 2}
\end{align}

We have to find the eigen vectors. For that we first have to find the eigen values using the characteristic equation

\begin{align}
det(\Vec{A} - \lambda \Vec{I}) = 0 \\
(1 - \lambda)(2 - \lambda) = 0 \\
\end{align}

Hence,
\begin{align}
\lambda_1 = 1, \quad \lambda_2 = 2
\end{align}

Now, to find the eigenvector for each eigen value;

\begin{align}
\Vec{A}\Vec{v} = \lambda\Vec{v} \\
\implies (\Vec{A} - \lambda I)\vec{v} &= 0 \\
\myvec{1-\lambda & 2 \\ 0 & 2 - \lambda}\Vec{v} = \myvec{0 \\ 0}
\end{align}

Writing as an augmented matrix and row-reducing:
\begin{align}
\myvec{1-\lambda & 2 & \vline & 0 \\ 0 & 2-\lambda & \vline & 0}
\end{align}

For $\lambda=1$ we have;
\begin{align}
\myvec{0 & 2 & \vline & 0 \\ 0 & 1 & \vline & 0}
\end{align}

On solving, we get $\Vec{v}=\myvec{x \\ 0}$. \\
Thus, the first eigen vector is $\myvec{1 \\ 0}$.

For $\lambda=2$ we have;
\begin{align}
\myvec{-1 & 2 & \vline & 0 \\ 0 & 0 & \vline & 0}
\end{align}

On solving this, we get $\Vec{v}=\myvec{x \\ \frac{x}{2}}$. \\
Thus, the second eigen vector is $\myvec{1 \\ \frac{1}{2}}$.

Thus, when we get $a=0$ and $b=\frac{1}{2}$. So $a+b=\frac{1}{2}$.
Therefore, option d is correct.

\end{document}
